\documentclass[12pt]{extreport}
\usepackage[a4paper,left=3cm,right=2cm,top=2cm,bottom=2cm,headheight=16pt,headsep=8mm,footskip=10mm,includeheadfoot,heightrounded]{geometry}
\usepackage{fontspec}
\setmainfont{Times New Roman}
\usepackage{setspace}
\setstretch{1.35}
\usepackage{titlesec}
\titleformat{\chapter}[display]{\bfseries\fontsize{17}{22}\selectfont}{\chaptername\ \thechapter}{0.45em}{}
\titlespacing*{\chapter}{0pt}{0pt}{1.0em}
\usepackage{fancyhdr}
\pagestyle{fancy}\fancyhf{}\fancyhead[C]{\thepage}\renewcommand{\headrulewidth}{0pt}
\begin{document}
\pagenumbering{roman}
\setcounter{page}{6}
\chapter*{Danh sách từ viết tắt (List of Abbreviations)}
\begin{tabular}{p{0.22\textwidth}p{0.68\textwidth}}
AAA & BBB \\
CCC & DDD \\
\end{tabular}
\clearpage
\chapter*{Danh mục thuật ngữ chuyên ngành}
\begin{center}
\begin{tabular}{p{0.31\textwidth}p{0.59\textwidth}}
T1 & Chiến lược sinh nội dung theo thứ tự: xây dựng topology nhiệm vụ trước, sau đó mới sinh chi tiết phòng. \\
T2 & Đồ thị mô tả các phòng/nút nhiệm vụ và quan hệ kết nối, điều kiện mở khóa trong dungeon. \\
T3 & Bố cục tile bên trong một phòng, thể hiện tường, cửa, vật cản, thực thể và đường đi. \\
T4 & Không gian biểu diễn nén dùng để mã hóa đặc trưng ngữ nghĩa của phòng trước khi sinh mẫu mới. \\
T5 & Mã biểu diễn rời rạc trong latent space, thường được học bởi VQ-VAE cho dữ liệu tile. \\
T6 & Mô hình diffusion sinh mẫu dưới điều kiện đầu vào (ví dụ: ngữ cảnh phòng, ràng buộc đồ thị). \\
T7 & Cơ chế đưa thông tin từ mission graph vào mô hình sinh để giữ nhất quán tiến trình toàn cục. \\
T8 & Ràng buộc ngữ nghĩa gameplay, ví dụ tính hợp lệ khóa -- chìa, boss room hoặc thứ tự mở cửa. \\
T9 & Thành phần hướng dẫn quá trình sinh dựa trên luật và tín hiệu logic để giảm mẫu vi phạm. \\
T10 & Bước hậu xử lý dùng luật ký hiệu để sửa các lỗi cấu trúc hoặc tiến trình còn sót lại. \\
\end{tabular}
\end{center}
\end{document}
